\documentclass[12pt]{article}
\usepackage[utf8]{inputenc}
\usepackage[T1]{fontenc}
\usepackage{hyperref}
\title{Poqomchi' Language Resources for K-12 Curriculum Development}
\author{Compiled by Claude Sonnet 5.5 (claude-sonnet-5-5) for Andr\'as Kornai}
\date{}
\begin{document}
\maketitle

\section{Introduction}
Poqomchi' (ISO 639-3 \texttt{poh}) is a Mayan language of Guatemala with roughly 90,000 speakers. Search found a pedagogical dictionary, an online talking dictionary and a New Testament, but no school textbooks or grammar.

\section{Primary Educational Resources}
\begin{itemize}
\item Academy of Mayan Languages of Guatemala (ALMG) bilingual Poqomchi'--Spanish \emph{Diccionario Pedag\'ogico}, released May 2025 (news report). \url{https://agn.gt/comunidad-linguistica-poqomchi-lanza-diccionario-pedagogico/}
\end{itemize}

\section{Linguistic and Reference Materials}
\begin{itemize}
\item Marcel Dobbels, \emph{Diccionario Poqomchi'--Castellano} (2003, PROASE; 856 pp.), Glottolog entry. \url{https://glottolog.org/resource/reference/id/104700}
\item Endangered Languages Project page for Poqomchi'. \url{https://endangeredlanguages.com/node/112918}
\end{itemize}

\section{Religious and Community Resources}
\begin{itemize}
\item The Bible in Poqomchi' (New Testament; Wycliffe), Davar Partners portal. \url{https://dbs.org/bibles/POHBSG}
\end{itemize}

\section{Digital and Audio Resources}
\begin{itemize}
\item Poqomchi' Tusq'orik, Living Dictionary (audio talking dictionary, about 1,845 entries with English and Spanish translations). \url{https://livingdictionaries.app/poqomchi/about}
\end{itemize}

\section{Curriculum Development Context}
The language is used by the whole community, and the ALMG dictionary is a natural base for vocabulary. No grammar or school materials were found, so these would need to be developed with ALMG and community teachers.

\end{document}